% PRISM-Hub 최종 보고서 조판 설정(XeLaTeX + xetexko)
\usepackage[a4paper,top=25mm,bottom=25mm,left=25mm,right=22mm,headheight=14pt]{geometry}
\usepackage{fontspec}
\usepackage{xetexko}
\setmainfont{Noto Serif CJK KR}[BoldFont={Noto Serif CJK KR Bold}]
\setsansfont{Noto Sans CJK KR}[BoldFont={Noto Sans CJK KR Bold}]
\setmonofont{DejaVu Sans Mono}[Scale=0.85]
\setmainhangulfont{Noto Serif CJK KR}[BoldFont={Noto Serif CJK KR Bold}]
\setsanshangulfont{Noto Sans CJK KR}[BoldFont={Noto Sans CJK KR Bold}]
\linespread{1.32}
\setlength{\parindent}{0pt}
\setlength{\parskip}{0.55em}

\usepackage{amsmath}
\usepackage{graphicx}
\usepackage{float}
\usepackage{array}
\usepackage{longtable}
\usepackage{booktabs}
\usepackage{ragged2e}
\usepackage{enumitem}
\setlist{itemsep=0.15em,topsep=0.2em,parsep=0.1em,leftmargin=1.6em}
\usepackage[table]{xcolor}
\definecolor{prismblue}{HTML}{1F4E79}
\definecolor{rowgray}{HTML}{F2F2F2}
\usepackage{caption}
\captionsetup{font=small,labelfont=bf,labelsep=period,justification=centering,singlelinecheck=false}
\captionsetup[table]{position=top,skip=4pt}
\captionsetup[figure]{position=bottom,skip=4pt}
\usepackage{titlesec}
\titleformat{\chapter}[hang]{\sffamily\bfseries\LARGE\color{prismblue}}{\thechapter.}{0.6em}{}
\titlespacing*{\chapter}{0pt}{0pt}{1.2em}
\titleformat{\section}[hang]{\sffamily\bfseries\large}{\thesection}{0.6em}{}
\titleformat{\subsection}[hang]{\sffamily\bfseries\normalsize}{\thesubsection}{0.6em}{}
\titlespacing*{\section}{0pt}{1.2em}{0.4em}
\titlespacing*{\subsection}{0pt}{0.9em}{0.3em}
\usepackage{fancyhdr}
\pagestyle{fancy}
\fancyhf{}
\fancyhead[R]{\small\sffamily\color{gray}PRISM-Hub 최종 보고서}
\fancyfoot[C]{\small\thepage}
\renewcommand{\headrulewidth}{0.3pt}
\fancypagestyle{plain}{\fancyhf{}\fancyfoot[C]{\small\thepage}\renewcommand{\headrulewidth}{0pt}}
\usepackage{tocloft}
\renewcommand{\cftchapfont}{\sffamily\bfseries}
\renewcommand{\cftchapaftersnum}{.}
\setlength{\cftchapnumwidth}{2.4em}
\setlength{\cftsecnumwidth}{3.0em}
\renewcommand{\contentsname}{목차}
\renewcommand{\listfigurename}{그림 목차}
\renewcommand{\listtablename}{표 목차}
\renewcommand{\bibname}{참고문헌}
\usepackage[hidelinks,unicode,bookmarksnumbered]{hyperref}

% 표·그림 번호: 장-순번(예: 표 2-1)
\renewcommand{\thetable}{\thechapter-\arabic{table}}
\renewcommand{\thefigure}{\thechapter-\arabic{figure}}
\captionsetup[table]{name=표}
\captionsetup[figure]{name=그림}

% 저장소 최상위(그림 경로용. 본문에는 쓰지 않는다)
\newcommand{\repo}{../../..}

% 표 아래 주·자료 줄
\newcommand{\srcnote}[1]{\par\vspace{-0.2em}{\footnotesize\raggedright #1\par}\vspace{0.6em}}
% 표 칸 정렬
\newcolumntype{L}[1]{>{\RaggedRight\arraybackslash}p{#1}}
\newcommand{\tabhead}{\rowcolor{prismblue!12}}
% 판정 강조
\newcommand{\pass}{\textbf{통과}}
\newcommand{\fail}{\textbf{불통과}}
\setlength{\LTpre}{0.4em}
\setlength{\LTpost}{0.2em}
\renewcommand{\arraystretch}{1.18}
\newfontfamily\fallbackfont{DejaVu Serif}
